\documentclass[11pt]{article}
\usepackage[letterpaper,margin=.68in,headheight=16pt,headsep=.18in,footskip=.28in]{geometry}
\usepackage{amsmath,array,enumitem,fancyhdr,hyperref}
\pagestyle{fancy}\fancyhf{}\renewcommand{\headrulewidth}{.3pt}
\fancyhead[L]{\small Bellingham Math Circle / Week 15 / Return visits}
\fancyhead[R]{\small Adult guide}
\fancyfoot[L]{\small RV-15-FAC-v1 / Draft and unpiloted}\fancyfoot[R]{\small\thepage}
\setlength{\emergencystretch}{2em}
\setlength{\parindent}{0pt}\setlength{\parskip}{7pt}
\setlist[itemize]{leftmargin=1.3em,itemsep=3pt,topsep=4pt}
\hypersetup{hidelinks,pdftitle={Week 15 return-visit facilitator guide}}
\begin{document}
{\Large\bfseries Mathematical overview}\par\medskip
\textbf{Facts and limits.} Taxi distance permits horizontal/vertical travel from any point, including inside a cell; fractional edges count as fractional steps. It uses \(|\Delta x|+|\Delta y|\), rather than a straight string. Between (0,0) and (2,2), ties occupy whole opposite quadrants, not only a line. The closed nearest region can fail convexity: (4,0) and (0,4) tie, while their midpoint (2,2) belongs strictly to the second site.

For greatest straight-line distance, each pairwise comparison is still a half-plane, so each farthest-site region is convex, but some sites have no region. For four square corners (\(\pm2,\pm2\)) plus center E, E is never farthest. Corner cells are opposite their corners. In contrast, every distinct site owns itself under nearest-distance rules.

In the bounded square with its four corner sites, the largest distance to the nearest site occurs uniquely at the center, distance \(\sqrt8\). Once that center is also a site, the new optimum is distance 2, attained only at the four edge midpoints. Exact quadrant arguments establish these facts; a sampled map alone cannot prove a continuum maximum.

\textbf{Children's destination.} Identify how the metric changes a map, see an empty cell under a reversed comparison, and optimize the worst nearby distance with a construction and a bound.

\textbf{Entry map.} Problem 1 (page 1): K--1 and up with string and oral ``farthest'' rules. Problem 2 (page 2): grades 2--3 and up, whole grid steps first, fractional lengths and midpoint later. Problems 3--4 (pages 3--4): grades 4--5, concrete candidate locations and a picture-based all-points bound; square roots and coordinate algebra are adult notation.

\textbf{New versus base.} Base pages make straight-line nearest cells, ties, insertions and prescribed regions. These visits change the distance, reverse nearest to farthest, and ask for maximum clearance.
\newpage {\Large\bfseries Prepare a return visit}\par\medskip
\textbf{Status and use.} Three investigations extend one familiar theme. These companions are separate from the base packets and may support several visits. All pages are new and unpiloted. Printed labels describe prerequisites, not age restrictions. Offer one investigation’s pages at a time; completing every page is not the goal.

\textbf{Materials and preparation.} For each pair, pencil/eraser, ruler, about 12 inches of nonstretch string, and spare grid/tracing paper. Reserve six sets for eleven children. All four printed frames are 4.5 inches square with equal x/y scale. Measurements are to dot centers. A printed frame bounds the clearance task but only displays part of the plane in the whole-region questions. Exact ties use symmetry; stock sizes and string procedure have not been rehearsed.

\textbf{Staffing.} Current context: eleven children, fixed KK11 / 3333 / 445 tables, one adult anchored to each table. Use pairs; the upper third child rotates as reader or checker. Routine adult reading, arrow drawing or recording can leave the mathematical decisions with children. If adults must choose every move, use the earlier entry rather than run the investigation for them.

\textbf{Short common launch.} Let children stretch strings between dots. For the taxi page use its printed non-task route: P=(0.5,0.5), Q=(3,1), two and a half steps across plus half a step up, total three. P starts inside a cell; fractions measure length. Show input, route and total in one unchanged reference. Then compare two string lengths from a trial point to sites and ask which is greatest. Announce the relevant comparison rule at each table/page; do not silently carry grid steps into a straight-line task.

\textbf{Suggested first route.} Begin with whichever physical comparison is familiar. Problem 1 is the youngest table entry. Grid-step classification requires keeping steps aligned with the grid. Save a proof of the whole-square clearance optimum for a return visit after genuine candidate testing.

\textbf{Flexible timing.} Allow 3--5 minutes of material play and 2--4 minutes for the common demonstration, then 15--25 minutes for one investigation and 5 minutes to share. A longer second visit can spend 30--45 minutes on the same investigation, including unfinished examples or its explanation. These are planning estimates, not observed timings. Do not require all three within one hour.

\textbf{Questions and hints.} Start with ``What can you try?'' and ``What would count as success?'' Check the actual rule before giving a strategy. Optional questions and hints are keyed below; use them only after a real attempt. A counter argument or drawing may be a complete explanation.

\textbf{Source and pilot record.} Mathematical examples and arguments here are original local follow-ups to the current base theme. Consulted teaching source: \emph{Math Circle by the Bay}, Preface pp. vii--x (local PDF pp. 8--11), on interaction, manipulatives, hints, explanations and themes spanning multiple years. Our exact launch, entry route and timing are design inferences. Year-1 \emph{Handouts 6}, Problems 6.2--6.6, provided a local example of connecting representations. Record actual problem, starting route, examples tried, what needed adult rescue, and what children want to resume. Distinguish observations from hypotheses. Counter fit, materials stock, procedures and classroom response have not been physically tested.
\newpage
\textbf{Problem 1: farthest ownership.} With square corners A=(-2,-2), B=(2,-2), C=(2,2), D=(-2,2), the farthest cells are: A owns x\(\geq\)0,y\(\geq\)0; B owns x\(\leq\)0,y\(\geq\)0; C owns x\(\leq\)0,y\(\leq\)0; D owns x\(\geq\)0,y\(\leq\)0. Axis boundaries are shared; all four tie at the center.

Every point is farthest from the corner opposite its quadrant. Moving the corner away horizontally increases the squared horizontal distance, and moving it away vertically increases the squared vertical distance. The two comparisons are independent.

\textbf{Why E has no cell.} Average the four squared corner distances from any point (x,y). The average is \(x^2+y^2+8\), while the squared distance to E=(0,0) is \(x^2+y^2\). At least one corner distance is therefore strictly larger than E's. This is a proof for every point in the plane. An elementary drawing argument can split into quadrants and compare E with the opposite corner, whose horizontal and vertical separations are both greater.

\textbf{If needed.} Ask where a point very near one corner should belong under ``farthest.'' Keep the rule card visibly reversed. String experiments suggest the opposite corner; the quadrant comparison explains the whole region.

\textbf{Return question.} Which added points can have empty farthest regions? Interior points of a convex hull cannot be uniquely farthest; do not claim the precise all-case classification without a new argument.
\newpage
\textbf{Problem 2: taxi ties.} For A=(0,0), B=(2,2), within the small square \([0,2]^2\), the grid distances are x+y and 4-x-y. Thus their tie line there is x+y=2. But outside that square, the same line description fails: for every x\(\geq\)2, y\(\leq\)0, both distances are x-y; the opposite quadrant x\(\leq\)0, y\(\geq\)2 is tied as well.

A precise whole-plane test is
\[
f(x)+f(y)\leq0\quad\hbox{for A's closed region},\qquad
f(t)=|t|-|t-2|.
\]
Here f is -2 below 0, rises linearly from -2 to 2 between 0 and 2, and stays 2 above 2. Equality means tie. On the printed frame \([-2,6]^2\), the 81 crossings comprise 15 strict A, 35 strict B and 31 tied AB. Sixteen entire unit cells are tied. This is adult compact notation; children can compare routes across the nine strips made by x=0,2 and y=0,2.

\textbf{The convexity surprise.} Both (4,0) and (0,4) have distance four to A and B. Their straight joining segment contains B=(2,2), where distance to B is zero and to A is four. Thus a closed nearest region under taxi distance need not contain the whole segment between two of its points.

\textbf{If needed.} Put a small token on each counted grid step, then collect the tokens to compare distances. Any shortest grid route has the same horizontal and vertical step counts; wandering routes do not define the distance.

\textbf{Extension.} Move the second site onto the same horizontal grid line. Its tie set returns to a vertical line. The large tie regions arise from the alignment of two changing coordinate contributions, not an imprecision in measurement.
\newpage
\textbf{Problems 3--4: the clearest place.} With only the four square corners as sites, choose the center. Its nearest-site distance is \(\sqrt{2^2+2^2}=\sqrt8\). Every point lies in a quadrant where its distances horizontally and vertically from that quadrant's corner are at most two. Its distance to that corner is therefore at most \(\sqrt8\); equality in both coordinates forces the center. This proves uniqueness and optimality.

Problem 4 adds a center site; the four side midpoints each have nearest-site distance two, to the center and the adjacent corners. To prove no better point exists, work in the upper-right quarter \(0\leq x,y\leq2\). If x+y\(\leq\)2, distance to the center is at most x+y, hence at most two. If x+y\(\geq\)2, the horizontal and vertical separations from the corner (2,2) sum to \(4-x-y\leq2\), so distance to that corner is at most two. Equality of Euclidean length with the sum of nonnegative horizontal/vertical lengths requires one separation to vanish. These conditions give only (2,0) and (0,2). Reflect into the other quadrants to obtain all four edge midpoints.

\textbf{Units and domain.} In the adult coordinate normalization the square has side four. On the printed 4.5-inch square, the first optimum distance is \(4.5/\sqrt2\approx3.182\) inches and the second is 2.25 inches. Only the candidate point must lie in or on the square; its clearance circle need not fit inside the frame.

\textbf{If needed.} Give children string circles centered on candidate points, or compare the nearest of five lengths. Suggest testing the center and side midpoints only after their own candidates. The proof is a return question; a good candidate search alone can be a satisfying first visit.

\textbf{Extension.} Where should a second new site go after choosing the center? This is a different optimization because the existing reference state has changed. Keep it explicit and allow multiple plausible experiments.

\textbf{Checks.} Exact arguments above settle the continuum claims; a 0.05-unit grid audit confirms the examples and finds the same optimum locations. Sampling is supporting evidence, not the proof.
\end{document}
